\chapter{Introduction}

\section{Background and Motivation}

The security of modern digital communications relies heavily on cryptographic algorithms that are based on mathematical problems considered computationally intractable for classical computers. Public-key cryptography systems such as RSA (Rivest-Shamir-Adleman) and Elliptic Curve Cryptography (ECC) have been the backbone of secure internet communications for decades, protecting everything from web browsing to financial transactions.

However, the emergence of quantum computing poses an existential threat to these cryptographic foundations. In 1994, Peter Shor developed a quantum algorithm capable of efficiently solving both the integer factorization problem and the discrete logarithm problem—the very mathematical assumptions underlying RSA and ECC \cite{Nielsen2010}. While large-scale quantum computers do not yet exist, rapid progress in quantum hardware development has made the threat increasingly tangible.

In response to this challenge, the cryptographic community has been developing post-quantum cryptography (PQC)—cryptographic algorithms designed to be secure against both classical and quantum computers. In 2016, the National Institute of Standards and Technology (NIST) initiated a standardization process to identify and standardize post-quantum cryptographic algorithms \cite{NIST-PQC}. After multiple rounds of evaluation, CRYSTALS-Kyber was selected in 2022 as the primary algorithm for public-key encryption and key establishment.

\section{Project Objectives}

The primary objectives of this graduation project are:

\begin{enumerate}
\item \textbf{Theoretical Understanding:} To gain comprehensive knowledge of cryptographic handshake protocols, classical public-key cryptography (specifically ECC), and post-quantum cryptography (focusing on CRYSTALS-Kyber).

\item \textbf{Mathematical Foundation:} To understand the mathematical underpinnings of both classical and post-quantum cryptographic systems, including elliptic curve operations and lattice-based cryptography fundamentals.

\item \textbf{Implementation:} To develop a working demonstration system that implements both ECC and Kyber-based key exchange mechanisms with proper client-server architecture.

\item \textbf{Performance Optimization:} To explore and implement SIMD (Single Instruction Multiple Data) optimizations using AVX2 instructions for accelerating Kyber's polynomial arithmetic operations.

\item \textbf{Comparative Analysis:} To benchmark and compare the performance characteristics of classical ECC versus post-quantum Kyber under different implementation strategies.

\item \textbf{Practical Demonstration:} To create an interactive visualization system demonstrating real-time cryptographic handshakes and their performance metrics.
\end{enumerate}

\section{Scope and Methodology}

This project encompasses three main phases:

\subsection{Research Phase}
We conducted extensive research into cryptographic theory, studying foundational texts and recent academic papers on both classical and post-quantum cryptography. Key topics included:
\begin{itemize}
\item Public-key cryptography fundamentals
\item Elliptic curve mathematics and the secp256k1 curve
\item Lattice-based cryptography and the Learning With Errors (LWE) problem
\item Module-LWE and the Kyber algorithm specification
\item NIST post-quantum standardization process
\end{itemize}

\subsection{Implementation Phase}
The implementation phase involved developing a complete cryptographic library featuring:
\begin{itemize}
\item Pure Python implementations for algorithm clarity and verification
\item Optimized C implementations with SIMD acceleration for performance
\item Client-server architecture for demonstrating key exchange
\item Proper protocol design for handshake mechanisms
\end{itemize}

\subsection{Evaluation Phase}
The final phase focused on testing, benchmarking, and creating visualization tools to demonstrate and compare the implemented algorithms.

\section{Report Organization}

The remainder of this report is organized as follows:

\textbf{Chapter 2} provides the theoretical background, covering cryptographic handshake protocols, elliptic curve cryptography, and the mathematical foundations of CRYSTALS-Kyber.

\textbf{Chapter 3} describes the implementation details, including system architecture, algorithm implementations, SIMD optimizations, and the client-server protocol design.

\textbf{Chapter 4} presents performance analysis and comparative benchmarks between ECC and Kyber implementations.

\textbf{Chapter 5} concludes the report with a summary of findings, contributions, and future work.